The learning-rate schedule is one of the cheapest knobs in neural-network training and one of the easiest to get wrong. Transformers were introduced with an explicit warmup-then-decay rule~\cite{vaswani2017attention}; linear warmup had earlier been recommended for large-minibatch SGD~\cite{goyal2017accurate}; and cosine annealing~\cite{loshchilov2016sgdr} combined with decoupled weight decay in AdamW~\cite{loshchilov2017decoupled,kingma2014adam} is today's default recipe. Why warmup helps is still debated: one explanation attributes it to the high variance of Adam's adaptive step early in training~\cite{liu2019variance}, others analyse how it lets a network tolerate a larger peak learning rate~\cite{kalra2024why} or how much of it GPT training can do without~\cite{kosson2024analyzing}. A recent alternative removes the decay schedule altogether: Schedule-Free AdamW~\cite{defazio2024road} replaces it by iterate averaging and is advertised as needing no stopping-time-dependent schedule.

Optimiser benchmarks repeatedly find that conclusions depend on the tuning protocol and schedule~\cite{schmidt2020descending,dahl2023benchmarking}. A practitioner's question that these large studies leave open is more modest: \emph{if I train a small model for a short, fixed budget and cannot afford to tune the peak learning rate carefully, which schedule loses least when the peak rate is wrong?} We answer it for a deliberately small, fully reproducible setting that runs on a laptop CPU: a two-layer character-level GPT on Tiny Shakespeare for \rhval{const/steps} steps.

\paragraph*{Contributions.} (i) A controlled comparison of four schedules, all sharing one hand-written AdamW, over five peak learning rates and \rhval{const/n_seeds} seeds each, with a defined sensitivity metric. (ii) Evidence that, in this setting, much of the apparent advantage of cosine decay over constant and step schedules at large peak rates is explained by warmup, which the baselines as originally specified lack; we separate the two with a warmup ablation. (iii) A negative result: Schedule-Free AdamW reaches the lowest loss but is not less sensitive to the peak LR than warmup+cosine.

\paragraph*{Hypotheses} (registered before the runs in \texttt{proposal.md}). \textbf{H1}: at each system's best main peak LR, warmup+cosine has lower validation loss than constant LR. \textbf{H2}: Schedule-Free AdamW has lower peak-LR sensitivity than warmup+cosine. \textbf{H3}: at the largest main peak LR ($10^{-2}$), removing warmup from cosine hurts.
